\section{Related Work}\label{sec:related}

The literature most relevant to this manuscript sits at the intersection of two strands that do not always talk to each other. One strand evaluates LEZs or related mobility restrictions as public policy and emphasizes treatment timing, comparison groups, and the interpretation of post-policy changes in regulated pollutants. Recent examples report measurable improvements in traffic-related air quality while also stressing trade-offs, spillovers, and heterogeneous local effects \cite{PRIETORODRIGUEZ2022237,SARMIENTO2023105014}. This strand is useful for motivating the policy question, but its identification logic often differs from the forecasting-based design used here.

The second strand focuses on air-quality forecasting itself. Here the dominant concern is predictive accuracy rather than policy evaluation. Prior work has explored machine-learning and deep-learning models for pollutant prediction under different feature sets, including meteorology-only inputs, traffic-related covariates, and mixed-source data \cite{SAMAD2023119987,Rahman2024}. These studies show that multivariate forecasting can capture meaningful pollutant dynamics, but they usually stop at prediction performance and do not translate the forecasts into a policy counterfactual.

This paper uses forecasting to test post-policy deviations rather than as an end in itself. Models are trained only on the pre-policy regime and then used to estimate what pollution levels would have looked like in the absence of the post-policy shift. In that sense, the forecasting model plays the role of a counterfactual generator. The approach is closer in spirit to event-based designs that compare observed outcomes with an expected baseline \cite{Guo2021}, but it remains prediction-driven rather than fully econometric. The gap addressed here is the relative scarcity of compact, robustness-oriented studies that combine both perspectives without letting the conclusion rest on a single preferred specification.
